\documentclass[11pt]{article}
\usepackage[T1]{fontenc}
\usepackage{textcomp}
\usepackage[margin=0.75in]{geometry}
\usepackage{amsmath,amssymb,amsthm}
\usepackage{array,booktabs}
\usepackage{hyperref}
\setlength{\parindent}{0pt}
\newtheorem{theorem}{Theorem}
\theoremstyle{definition}
\newtheorem{definition}{Definition}
\title{Flow Proof Artifact: Nat-order-trichotomy}
\author{Generated by \texttt{flow doc proof}}
\date{Flow Proof Book}
\begin{document}
\maketitle
Any two naturals are comparable by less-or-equal in one direction.\\[0.5em]
\textbf{Source.} peano — https://en.wikipedia.org/wiki/Trichotomy\_(mathematics)\\[0.5em]
\bigskip
\noindent{\large \textbf{Derived fact 1.} \textit{For all a and b, either a ≤ b or b ≤ a} \hfill the natural numbers \cdot order \cdot one side is less than or equal to the other}\par
\smallskip\noindent\textit{Coordinate: the natural numbers \cdot order \cdot one side is less than or equal to the other}\par
\smallskip\noindent\textit{Source: peano/induction — Gries \& Schneider, Ch. 3}\par
\smallskip\noindent\textit{Built on: less-or-equal is reflexive, for order on the natural numbers, less-or-equal is transitive, for order on the natural numbers, every number is below its successor, for order on the natural numbers}\par
\medskip
\noindent\textbf{Goal.} For all a and b, either a ≤ b or b ≤ a\par
\noindent$\displaystyle \forall a \in \mathbb{N} \forall b \in \mathbb{N}\quad b \le a$\par
\medskip
\noindent
\begin{tabular}{@{} >{\raggedright\arraybackslash}p{0.06\textwidth} >{\raggedright\arraybackslash}p{0.47\textwidth} >{\raggedleft\arraybackslash}p{0.41\textwidth} @{}}
\textbf{\#} & \textbf{Proof} & \textbf{Mathematics} \\
\midrule
\textbf{1.} & We prove that one side is less than or equal to the other for order on the natural numbers. & \\[0.45em]
\textbf{2.} & We split into exhaustive cases — the claim must hold in each one. & \\[0.45em]
\textbf{3.} & Case 1 (see step 2): suppose a  equals  b. & \\[0.45em]
\textbf{4.} & We invoke the derived fact governing order on the natural numbers: less-or-equal is reflexive, for order on the natural numbers (instantiated for a). & \\[0.45em]
\textbf{5.} & From step 3 and step 4, this implies a is at most b in this case. & $\displaystyle a \le b$ \\[0.45em]
\textbf{6.} & Case 2 (see step 2): suppose a < b. & \\[0.45em]
\textbf{7.} & From step 6, this implies a is at most b in this case. & $\displaystyle a \le b$ \\[0.45em]
\textbf{8.} & Case 3 (see step 2): neither disjunct holds. & \\[0.45em]
\textbf{9.} & We invoke the derived fact governing order on the natural numbers: less-or-equal is reflexive, for order on the natural numbers (instantiated for b). & \\[0.45em]
\textbf{10.} & From step 8 and step 9, this implies b is at most a. Together with the other cases (step 3, step 6, and step 8), the goal is discharged. Hence proven. & $\displaystyle b \le a$ \\[0.45em]
\bottomrule
\end{tabular}
\label{thm:Nat-order-one_side_is_less_than_or_equal_to_the_other}
\par\medskip\hrule
\end{document}
